\scanpage{106}
de $\unsure{\xi|\zeta}$ correspond au choix d'un $x\in X$ et d'une ext. sép. close \deleted{$\bar K$} \inserted{$\Omega$} de $k(x)$ (pas néc. \deleted{\unsure{une}} clôture alg. de $k(x)$) et on note alors aussi $\mathcal E_X^\Omega$ au lieu de $\mathcal E_X^{\unsure{\xi|\zeta}}$ ($\Omega$ sous-entendu en tant qu'extension de $k(x)$, donc avec sa structure de $k(x)$ algèbre). On aura \inserted{un isom can.}
\[
(4)\qquad \mathcal E_{\unsure{\alpha|X}}^\Omega\simeq\mathcal E_X^{\overline{k(x)}},
\]
où $\overline{k(x)}$ est la clôture alg. de $k(x)$ dans $\Omega$. Bien sûr, si $X=\operatorname{Spec}(A)$, on notera aussi $\mathcal E_A^\Omega$ --- notation surtout utilisée pour $\mathcal E_K^{\bar K}$, $K$ un corps, $\bar K$ une clôture alg. ou sép. de $K$.

Si $X$ est un $Y$-schéma, \inserted{$Y$ connexe,} on notera
\[
(5)\qquad \mathcal E(f):\mathcal E_X\mathrel{\text{---}}\mathcal E_Y\qquad\text{où }f:X\longrightarrow Y
\]
($\mathcal E_X$ est un foncteur en $X$
\[
(\text{Sch conn})\mathrel{\text{---}}\unsure{\text{grpes ext}}
\]
\margin{105}
